\clearpage
\fancyhead[L]{\small MONDSTERNE.}
\hypertarget{native-448-mondsterne}{}
\begin{center}
{\large MONDSTERNE.}
\end{center}

\begin{center}
Astronomische Nachrichten. Band II. Nr. 26. Seite 17..20. 1823 Februar.
\end{center}

\begin{center}
Berechnung der Meridiandifferenz zweier Orte aus correspondirenden Mondsculminationen.\\
Von Herrn Professor \textsc{Nicolai}, Director der Sternwarte in Mannheim [an \textsc{Schumacher.}].
\end{center}

Ihrer gütigen Aufforderung, Ihnen die Art mitzutheilen, wie ich aus den correspondirenden
Mondsculminationen die Meridiandifferenz berechne, will ich heute durch diese Zeilen
nachzukommen suchen.

Diese Methode beruht, wie bekannt, darauf, dass an zwei Orten der wahre
Rectascensions-Unterschied des einen Mondsrandes und einiger vorher verabredeter Fixsterne
beobachtet wird. Es sei der Unterschied der wahren geraden Aufsteigung des Mondsrandes und
eines Fixsterns (positiv angenommen, wenn die Rectascension des Monds grösser ist als die des
Sterns) an dem westlichen Beobachtungsorte $=t$ an dem östlichen $=\tau$, so wird $t-\tau$ der
wahre Rectascensions-Unterschied des Mondsrandes für die zwischen beiden Beobachtungen
verflossene Zeit sein. Sind mehrere Fixsterne an beiden Orten verglichen worden, so wird man
aus den verschiedenen $t-\tau$ das arithmetische Mittel nehmen in so fern sie an beiden Orten
mit gleicher Genauigkeit beobachtet worden sind. Wie man zu verfahren hat, wenn dieses
rücksichtlich der Zahl der beobachteten Fäden nicht der Fall ist, werden wir weiter unten
sehen. --- Offenbar kommt es nun aber darauf an, den wahren Rectascensions-Unterschied des
\emph{Mittelpunkts} des Mondes zu haben, welcher von dem vorhin erhaltenen um etwas
verschieden sein wird, indem der Halbmesser des Mondes in AR. in der Zwischenzeit der
Beobachtungen sich ändert, theils durch die Aenderung des wahren Mondshalbmessers selbst,
theils durch die Aenderung der Declination des Mondes. Bezeichnet man für die Culminationszeit
des Mondes an dem westlichen Beobachtungsorte den \emph{wahren} (vom Mittelpunkt der Erde aus
gesehenen) Halbmesser desselben mit $r$, seine wahre Declination mit $d$, und für die
Culminationszeit an dem östlichen Orte eben diese Grösse mit $\rho$ und $\delta$, so wird

\[
t-\tau \mp \tfrac{1}{15}\left(\frac{r}{\cos d}-\frac{\rho}{\cos\delta}\right) = \Delta
\]
\noindent der wahre Rectascensionsunterschied des Mondsmittelpunkts für die beiden
Beobachtungsmomente sein. Das untere Zeichen in diesem Ausdrucke gilt, wenn der westliche
oder erste, das obere, wenn der östliche oder zweite Mondsrand beobachtet ist. In allen den
Fällen, wo die Meridiandifferenz der beiden Beobachtungsorte nicht sehr gross ist, wird man
diese Correction ganz vernachlässigen können; auch
\SourceLeaf{449}%
sieht man leicht, dass man die absoluten Werthe von $r$ und $d$ nicht in grosser Schärfe zu
kennen braucht, da es hier nur auf die \emph{Aenderung} dieser Grössen in der Zwischenzeit der
Beobachtungen ankommt.

Um aus $\Delta$ den Meridianunterschied der beiden Beobachtungsorte zu erhalten, bedürfen wir
nur eines einzigen Elementes, nemlich der stündlichen Bewegung des Mondes in AR. Bezeichnet
man diese durch $h$, die Meridiandifferenz in Zeit durch $x$, so ist

\[
x = \left(\frac{1}{h}\,m\cdot 15\cdot 3600 - 1\right)\Delta = n\Delta
\]
\noindent wenn man den Factor $\frac{1}{h}m\cdot15\cdot3600-1=n$ setzt. In diesem Ausdrucke hängt der
Factor $m$ von derjenigen Zeitart ab, für welche $h$ berechnet ist. Bedeutet nemlich $h$ die
Bewegung für eine Sternzeitstunde, so ist $m=1$; ist $h$ die Bewegung für eine mittlere
Sonnenstunde, so bezeichnet $m$ das Verhältniss des mittleren Sonnentages zum Sterntage, oder
ist $=1.00274$; ist endlich $h$ die Bewegung für eine wahre Sonnenstunde, so bedeutet $m$ das
jedesmalige Verhältniss des wahren Sonnentages zum Sterntage. --- Was nun die Berechnung der
stündlichen Bewegung $h$ betrifft, so ist hier zu bemerken, dass es nicht vortheilhaft sein
würde, wenn man sie aus den Mondstafeln durch unmittelbare Berechnung zweier Mondsörter,
welche nur eine Stunde von einander abstehen, herleiten wollte. Denn obgleich die Mondstafeln
auf Zehntheile von Secunden berechnet sind, so kann man doch, wegen der grossen Anzahl der zu
summirenden Gleichungen, einen aus ihnen berechneten Mondsort nie bis auf eine Secunde
verbürgen, und es ist demnach klar, dass die auf diese Art berechnete stündliche Bewegung
öfters merklich fehlerhaft erhalten werden würde. Weit sicherer, und in der That mit grosser
Schärfe, verfährt man daher, wenn man mehrere Mondsörter in \emph{beträchtlichen}
Zwischenräumen, z.\,B. von 12 zu 12 Stunden, berechnet, und aus diesen, vermittelst einer
\emph{scharfen Interpolation}, die stündliche Bewegung ableitet, wodurch sich die bei den
einzelnen Oertern noch statt findenden kleinen Fehler gleichmässig vertheilen werden. Zu
dieser Operation kann man sich nun ohne Bedenken der in den Ephemeriden von 12 zu 12 Stunden
berechneten Mondsörter bedienen. --- Bei der \emph{Veränderlichkeit} der stündlichen Bewegung
selbst, ist nun noch zu berücksichtigen, dass man das zu der Bestimmung der Meridiandifferenz
anzuwendende $h$ \emph{für das Mittel der beiden auf einerlei Meridian reducirten
Beobachtungszeiten} berechne. Es sei $A$ die gerade Aufsteigung des Mondes für die Stunde
$\theta$ unter einem gewissen Meridian, $A'$ dieselbe für die Stunde $\theta+1$ unter dem
nemlichen Meridian, so wird $A'-A$ oder $h$ die stündliche Bewegung des Mondes in
Rectascension für die Zeit $\theta+30^{\rm m}$ sein. Man wird also für den Meridian der Tafeln
oder Ephemeriden mehrere auf einander folgende stündliche Bewegungen berechnen, die beiden
Beobachtungszeiten vermittelst des beiläufig bekannten Meridianunterschiedes auf diesen
Meridian der Tafeln oder Ephemeriden reduciren, und für das Mittel der letztern aus jenen das
anzuwendende $h$ durch Interpolation erhalten.

Da die Vergleichungssterne sowohl als der Mond nicht an jedem Abend und von jedem Beobachter
immer an allen Fäden des Mittagsfernrohres beobachtet werden können, so ist klar, dass nicht
allen Resultaten eine gleiche Genauigkeit beigemessen werden kann, welche daher mit zu
berücksichtigen ist. Die Ausdrücke für die relative Genauigkeit der einzelnen Resultate wurden
mir früher schon von Herrn Hofrath \textsc{Gauss} gütigst mitgetheilt; sie ergeben sich leicht aus der
Methode der kleinsten Quadrate. Es sei der Mond an einem Orte an $l$, am andern an $l'$ ferner
seien die verschiedenen Sterne an einem Orte an $a, b, c$ etc. am andern an $a', b', c'$ etc.
Fäden beobachtet. Man setze

\[
\frac{l\,l'}{l+l'}=\lambda,\qquad \frac{a\,a'}{a+a'}=\alpha,\qquad \frac{b\,b'}{b+b'}=\ObservedBeta,\qquad \frac{c\,c'}{c+c'}=\gamma,
\]
\noindent etc. Nun hat man das Mittel der Resultate aus den einzelnen Vergleichungssternen so zu nehmen,
dass man ihnen resp. das Gewicht (Quadrat der Genauigkeit) $\alpha, \ObservedBeta, \gamma$ etc.
beilegt; das so erhaltene Resultat
\SourceLeaf{450}%
für $t-\tau$ oder für $\Delta$ hat alsdann das Gewicht $\dfrac{\sigma\lambda}{\sigma+\lambda}$, wenn man
$\alpha+\ObservedBeta+\gamma+\text{etc.}=\sigma$ setzt. Bezeichnet man nun den Coëfficienten, womit $\Delta$
multiplicirt wird, um die Meridiandifferenz zu erhalten, wie oben, mit $n$, so ist das Gewicht des
Resultates für diese letztere $=\dfrac{\sigma\lambda}{(\sigma+\lambda)nn}$, und das Gewicht des
Resultates aus allen Bestimmungen $=\Sigma\cdot\dfrac{\sigma\lambda}{(\sigma+\lambda)nn}$, oder der
wahrscheinliche Fehler des Endresultates, wenn man den wahrscheinlichen Fehler eines
Fadenantrittes $=\varepsilon$ setzt, ist

\[
=\frac{\varepsilon}{\BareRoot{\Sigma\,\dfrac{\sigma\lambda}{(\sigma+\lambda)nn}}}
\]

Wie man sieht, ist bei dieser Berechnung des wahrscheinlichen Fehlers vorausgesetzt, dass man allen
beobachteten Antritten gleiche Zuverlässigkeit beilegt, ohne weder Mond noch Sterne, noch die
verschiedenen Declinationen, noch die Beobachter zu unterscheiden. Es würde keine Schwierigkeit
haben, auf alle diese Umstände Rücksicht zu nehmen, wenn man die betreffenden Coëfficienten kennte.
Doch möchte alles dieses wol unerheblich sein. Was die verschiedenen Declinationen anbelangt, so
kann man ohne Bedenken alle im Zodiacus des Mondes beobachtete Antritte als gleich genau betrachten,
zumal, da die jedesmaligen Nebenumstände der Beobachtung, den Zustand der Luft, der tiefere oder
höhere Stand des Mondes und der Sterne u.\,s.\,w., hier weit mehr auf die Genauigkeit der Antritte
influiren werden, als die grösseren oder kleineren Declinationen. Uebrigens wird es zweckmässig sein,
für $s$ einen \emph{mittleren} Werth aus allen im Zodiacus des Mondes beobachteten Antritten in
obiger Endformel zum Grunde zu legen.

\begin{center}
{\small ---\,---\,---\,---\,---\,---\,---\,---\,---\,---\,---\,---\,---\,---\,---\,---\,---\,---\,---\,---\,---\,---}
\end{center}

Diese ganze Rechnung ist, wie Sie sehen, äusserst einfach, und auch, seitdem uns die Conn.\ de Tems
die Mühe der eignen Berechnung der Mondsrectascensionen erspart, sehr bequem und kurz. ---
Schliesslich bemerke ich noch, dass der Werth von $x+\Delta$ in \emph{vollkommener Strenge} aus
$\Delta$ eigentlich durch eine Reihe erhalten wird, die nach den ungeraden Potenzen von $\Delta$
fortläuft, und von welcher $\dfrac{1}{h}m\cdot15\cdot3600\Delta$ ($h$ für das Mittel der beiden
Beobachtungszeiten genommen) das erste Glied ist; indess ist das zweite von $\Delta^3$ abhängige
Glied schon so gering, dass es erst bei einer Meridiandifferenz von zwei bis dritthalb Stunden etwa
$0^{\mathrm{s}}1$ beträgt. 
Bei einem \emph{sehr grossen} Längenunterschiede würde ich aber ein \emph{indirectes} Verfahren
vorziehen, nemlich mit einer \emph{supponirten} Meridiandifferenz den ihr entsprechenden Unterschied
der Mondsrectascensionen berechnen, diesen berechneten mit dem wirklich beobachteten vergleichen, und
aus der Abweichung beider von einander, mittelst der stündlichen Bewegung, die an jene supponirte
Meridiandifferenz noch anzubringende Correction bestimmen.

\begin{flushright}
\textsc{Nicolai.}
\end{flushright}
